\documentclass[aspectratio=169]{beamer}
\usetheme{Madrid}
\usecolortheme{seahorse}

\usepackage{tikz}
\usepackage{xcolor}
\usepackage{listings}
\usepackage{amsmath}
\usepackage{booktabs}

% Custom colors
\definecolor{codegreen}{rgb}{0,0.6,0}
\definecolor{codegray}{rgb}{0.5,0.5,0.5}
\definecolor{codepurple}{rgb}{0.58,0,0.82}
\definecolor{backcolour}{rgb}{0.95,0.95,0.92}

% Code listing style
\lstdefinestyle{Rstyle}{
    backgroundcolor=\color{backcolour},   
    commentstyle=\color{codegreen},
    keywordstyle=\color{codepurple},
    numberstyle=\tiny\color{codegray},
    stringstyle=\color{codegreen},
    basicstyle=\ttfamily\footnotesize,
    breakatwhitespace=false,         
    breaklines=true,                 
    captionpos=b,                    
    keepspaces=true,                 
    numbers=left,                    
    numbersep=5pt,                  
    showspaces=false,                
    showstringspaces=false,
    showtabs=false,                  
    tabsize=2
}
\lstset{style=Rstyle}

\title{DirichletForestParallel}
\subtitle{High-Performance Compositional Data Regression}
\author{Your Name}
\date{\today}

\begin{document}

\begin{frame}
\titlepage
\end{frame}

% ============================================================================
% SLIDE 1: Overview
% ============================================================================
\begin{frame}{What is DirichletForestParallel?}

\begin{block}{Purpose}
Random forest regression for \textbf{compositional data} (e.g., proportions that sum to 1)
\end{block}

\vspace{0.3cm}

\textbf{Key Features:}
\begin{itemize}
\item Dirichlet distribution-based splits (likelihood maximization)
\item Parallel processing (Unix/Mac fork, Windows cluster)
\item Two prediction modes: Fast vs. Distributional
\item Multiple parameter estimation methods (MoM, MLE)
\item Weight-based interpretability
\end{itemize}

\vspace{0.3cm}

\textbf{Applications:}
\begin{itemize}
\item Microbiome composition analysis
\item Market share prediction
\item Budget allocation forecasting
\item Any data where components sum to 1
\end{itemize}

\end{frame}

% ============================================================================
% SLIDE 2: Input Data Structure
% ============================================================================
\begin{frame}[fragile]{Input Data Structure}

\begin{columns}[T]

\begin{column}{0.48\textwidth}
\textbf{Predictor Matrix (X)}

\begin{itemize}
\item Numeric matrix: $n \times p$
\item $n$ = observations
\item $p$ = features/predictors
\item Can be any real numbers
\end{itemize}

\vspace{0.2cm}

\begin{lstlisting}[language=R, basicstyle=\tiny\ttfamily]
# Example: 100 samples, 4 features
X <- matrix(rnorm(100 * 4), 
            nrow = 100, 
            ncol = 4)
\end{lstlisting}

\end{column}

\begin{column}{0.48\textwidth}
\textbf{Response Matrix (Y)}

\begin{itemize}
\item Numeric matrix: $n \times k$
\item $k$ = number of components
\item \textbf{Constraint:} Each row sums to 1
\item All values in $(0, 1)$
\end{itemize}

\vspace{0.2cm}

\begin{lstlisting}[language=R, basicstyle=\tiny\ttfamily]
# Example: 3 components (k=3)
library(MCMCpack)
Y <- rdirichlet(100, 
                alpha = c(2, 3, 4))

# Verify constraint
rowSums(Y)  # All equal to 1
\end{lstlisting}

\end{column}

\end{columns}

\vspace{0.4cm}

\begin{alertblock}{Important}
Y must be compositional: $\sum_{j=1}^{k} Y_{ij} = 1$ for all $i$, and $0 < Y_{ij} < 1$
\end{alertblock}

\end{frame}

% ============================================================================
% SLIDE 3: Key Parameters
% ============================================================================
\begin{frame}[fragile]{Key Function Parameters}

\begin{table}
\small
\begin{tabular}{lp{3cm}p{5cm}}
\toprule
\textbf{Parameter} & \textbf{Default} & \textbf{Description} \\
\midrule
\texttt{B} & 100 & Number of trees in forest \\
\texttt{d\_max} & 10 & Maximum tree depth \\
\texttt{n\_min} & 5 & Min samples per leaf node \\
\texttt{m\_try} & $\sqrt{p}$ & Features to try at each split \\
\texttt{method} & "mom" & Parameter estimation: "mom" or "mle" \\
\texttt{store\_samples} & FALSE & Storage mode (affects memory) \\
\texttt{n\_cores} & -1 & Parallel cores (-1 = all available) \\
\texttt{use\_leaf\_predictions} & TRUE & Use pre-computed vs. weight-based \\
\bottomrule
\end{tabular}
\end{table}

\vspace{0.3cm}

\begin{lstlisting}[language=R, basicstyle=\tiny\ttfamily]
# Example call
model <- DirichletForest_distributed(
  X, Y, 
  B = 100,           # 100 trees
  d_max = 10,        # Max depth 10
  n_min = 5,         # Min 5 samples per leaf
  method = "mom",    # Method of moments
  store_samples = TRUE,  # Enable weight analysis
  n_cores = 4        # Use 4 cores
)
\end{lstlisting}

\end{frame}

% ============================================================================
% SLIDE 4: Storage Modes
% ============================================================================
\begin{frame}{Storage Modes: Fast vs. Distributional}

\begin{columns}[T]

\begin{column}{0.48\textwidth}
\begin{block}{Fast Mode}
\texttt{store\_samples = FALSE}
\end{block}

\textbf{What's Stored:}
\begin{itemize}
\item Pre-computed $\alpha$ parameters
\item Pre-computed mean predictions
\item NO sample indices
\end{itemize}

\vspace{0.2cm}

\textbf{Pros:}
\begin{itemize}
\item 10-50× less memory
\item Instant predictions
\item Production-ready
\end{itemize}

\vspace{0.2cm}

\textbf{Cons:}
\begin{itemize}
\item Cannot analyze sample weights
\item Limited interpretability
\end{itemize}

\end{column}

\begin{column}{0.48\textwidth}
\begin{block}{Distributional Mode}
\texttt{store\_samples = TRUE}
\end{block}

\textbf{What's Stored:}
\begin{itemize}
\item $\alpha$ parameters + means
\item Sample indices in each leaf
\item Full training data (X, Y)
\end{itemize}

\vspace{0.2cm}

\textbf{Pros:}
\begin{itemize}
\item Weight-based predictions
\item Full interpretability
\item Uncertainty quantification
\end{itemize}

\vspace{0.2cm}

\textbf{Cons:}
\begin{itemize}
\item Higher memory usage
\item Slightly slower predictions
\end{itemize}

\end{column}

\end{columns}

\end{frame}

% ============================================================================
% SLIDE 5: Parameter Estimation Methods
% ============================================================================
\begin{frame}{Parameter Estimation Methods}

\textbf{Goal:} Estimate Dirichlet parameters $\boldsymbol{\alpha} = (\alpha_1, \ldots, \alpha_k)$

\vspace{0.3cm}

\begin{columns}[T]

\begin{column}{0.48\textwidth}
\begin{block}{Method of Moments (MoM)}
\texttt{method = "mom"}
\end{block}

\textbf{How it works:}
\begin{enumerate}
\item Calculate sample mean $\bar{y}_j$ and variance $s_j^2$
\item Estimate concentration: 
$$v = \frac{\bar{y}_1(1-\bar{y}_1)}{s_1^2} - 1$$
\item Set $\alpha_j = v \cdot \bar{y}_j$
\end{enumerate}

\vspace{0.2cm}

\textbf{Characteristics:}
\begin{itemize}
\item Very fast (closed-form)
\item Good for large datasets
\item May be less accurate
\end{itemize}

\end{column}

\begin{column}{0.48\textwidth}
\begin{block}{Maximum Likelihood (MLE)}
\texttt{method = "mle"}
\end{block}

\textbf{How it works:}
\begin{enumerate}
\item Initialize with MoM
\item Newton-Raphson optimization
\item Maximize log-likelihood:
$$\mathcal{L}(\boldsymbol{\alpha}) = \sum_i \log p(Y_i | \boldsymbol{\alpha})$$
\end{enumerate}

\vspace{0.2cm}

\textbf{Characteristics:}
\begin{itemize}
\item More accurate
\item Computationally expensive
\item Better for small datasets
\end{itemize}

\end{column}

\end{columns}

\end{frame}

% ============================================================================
% SLIDE 6: Parallel Processing
% ============================================================================
\begin{frame}[fragile]{Parallel Processing Architecture}

\textbf{Automatic platform detection:}

\vspace{0.3cm}

\begin{columns}[T]

\begin{column}{0.48\textwidth}
\begin{block}{Unix/Mac: Fork-based}
\end{block}

\begin{itemize}
\item Uses \texttt{mclapply}
\item Memory-efficient (copy-on-write)
\item Automatic setup
\item No cleanup needed
\end{itemize}

\vspace{0.2cm}

\begin{lstlisting}[language=R, basicstyle=\tiny\ttfamily]
# Unix/Mac automatically uses fork
model <- DirichletForest_distributed(
  X, Y, 
  B = 100,
  n_cores = 4  # Fork workers
)
\end{lstlisting}

\end{column}

\begin{column}{0.48\textwidth}
\begin{block}{Windows: PSOCK Cluster}
\end{block}

\begin{itemize}
\item Uses persistent cluster
\item Explicit worker setup
\item Higher overhead
\item Requires cleanup
\end{itemize}

\vspace{0.2cm}

\begin{lstlisting}[language=R, basicstyle=\tiny\ttfamily]
# Windows uses PSOCK cluster
model <- DirichletForest_distributed(
  X, Y, 
  B = 100,
  n_cores = 4
)

# IMPORTANT: Always cleanup!
cleanup_distributed_forest(model)
\end{lstlisting}

\end{column}

\end{columns}

\vspace{0.3cm}

\begin{alertblock}{Best Practice}
Set \texttt{n\_cores = -1} to auto-detect: uses \texttt{detectCores() - 1}
\end{alertblock}

\end{frame}

% ============================================================================
% SLIDE 7: Model Output Structure
% ============================================================================
\begin{frame}[fragile]{Model Output Structure}

\textbf{What you get after fitting:}

\vspace{0.3cm}

\begin{lstlisting}[language=R, basicstyle=\tiny\ttfamily]
model <- DirichletForest_distributed(X, Y, B = 100, store_samples = TRUE)
\end{lstlisting}

\vspace{0.3cm}

\begin{columns}[T]

\begin{column}{0.48\textwidth}
\textbf{Model Components:}
\begin{itemize}
\item \texttt{\$type} - "fork", "cluster", or "sequential"
\item \texttt{\$n\_cores} - Number of workers
\item \texttt{\$total\_trees} - Total trees in forest
\item \texttt{\$store\_samples} - Storage mode
\item \texttt{\$worker\_forests} - Individual worker forests (fork)
\item \texttt{\$cluster} - Cluster object (Windows)
\end{itemize}
\end{column}

\begin{column}{0.48\textwidth}
\textbf{Fitted Values (\texttt{\$fitted}):}
\begin{itemize}
\item \texttt{\$alpha\_hat} - Estimated $\alpha$ parameters
\item \texttt{\$mean\_based} - Direct mean predictions
\item \texttt{\$param\_based} - Normalized $\alpha$ predictions
\end{itemize}

\vspace{0.2cm}

\textbf{Residuals (\texttt{\$residuals}):}
\begin{itemize}
\item \texttt{\$mean\_based} - $Y - \hat{Y}_{\text{mean}}$
\item \texttt{\$param\_based} - $Y - \hat{Y}_{\text{param}}$
\end{itemize}
\end{column}

\end{columns}

\end{frame}

% ============================================================================
% SLIDE 8: Three Types of Predictions
% ============================================================================
\begin{frame}[fragile]{Three Types of Predictions}

\begin{block}{1. Alpha Parameters (\texttt{alpha\_predictions})}
Estimated Dirichlet parameters $\hat{\boldsymbol{\alpha}} = (\hat{\alpha}_1, \ldots, \hat{\alpha}_k)$
\end{block}

\begin{lstlisting}[language=R, basicstyle=\tiny\ttfamily]
pred <- predict_distributed_forest(model, X_test)
alpha_hat <- pred$alpha_predictions  # n_test x k matrix
\end{lstlisting}

\vspace{0.2cm}

\begin{block}{2. Mean-Based Predictions (\texttt{mean\_predictions})}
Direct average of compositional values: $\hat{Y}_{\text{mean}} = \frac{1}{B}\sum_{b=1}^{B} \bar{Y}_b$
\end{block}

\begin{lstlisting}[language=R, basicstyle=\tiny\ttfamily]
mean_pred <- pred$mean_predictions  # Sum to 1 by construction
\end{lstlisting}

\vspace{0.2cm}

\begin{block}{3. Parameter-Based Predictions}
Normalized alpha parameters: $\hat{Y}_{\text{param}} = \frac{\hat{\boldsymbol{\alpha}}}{\sum_j \hat{\alpha}_j}$
\end{block}

\begin{lstlisting}[language=R, basicstyle=\tiny\ttfamily]
param_pred <- pred$alpha_predictions / rowSums(pred$alpha_predictions)
\end{lstlisting}

\end{frame}

% ============================================================================
% SLIDE 9: Prediction Modes
% ============================================================================
\begin{frame}[fragile]{Prediction Modes (when store\_samples = TRUE)}

\begin{columns}[T]

\begin{column}{0.48\textwidth}
\begin{block}{Pre-computed Leaf Predictions}
\texttt{use\_leaf\_predictions = TRUE}
\end{block}

\textbf{How it works:}
\begin{enumerate}
\item Traverse tree to leaf
\item Return pre-computed $\alpha$ and mean
\item Average across all trees
\end{enumerate}

\vspace{0.2cm}

\textbf{Characteristics:}
\begin{itemize}
\item Very fast
\item Uses stored leaf values
\item Same as \texttt{store\_samples = FALSE}
\end{itemize}

\vspace{0.2cm}

\begin{lstlisting}[language=R, basicstyle=\tiny\ttfamily]
pred <- predict_distributed_forest(
  model, X_test,
  use_leaf_predictions = TRUE
)
\end{lstlisting}

\end{column}

\begin{column}{0.48\textwidth}
\begin{block}{Weight-Based Predictions}
\texttt{use\_leaf\_predictions = FALSE}
\end{block}

\textbf{How it works:}
\begin{enumerate}
\item Find all training samples in same leafs
\item Calculate normalized weights
\item Weighted average of training Y
\item Re-estimate $\alpha$ from weighted samples
\end{enumerate}

\vspace{0.2cm}

\textbf{Characteristics:}
\begin{itemize}
\item More accurate (full distribution)
\item Enables interpretability
\item Computationally expensive
\end{itemize}

\vspace{0.2cm}

\begin{lstlisting}[language=R, basicstyle=\tiny\ttfamily]
pred <- predict_distributed_forest(
  model, X_test,
  use_leaf_predictions = FALSE
)
\end{lstlisting}

\end{column}

\end{columns}

\end{frame}

% ============================================================================
% SLIDE 10: Weight Matrix Analysis
% ============================================================================
\begin{frame}[fragile]{Weight Matrix: Interpretability Tool}

\textbf{Only available with \texttt{store\_samples = TRUE}}

\vspace{0.3cm}

\begin{block}{What is the Weight Matrix?}
For $m$ test samples and $n$ training samples: $W_{m \times n}$ where $W_{ij}$ = influence of training sample $j$ on test sample $i$
\end{block}

\vspace{0.3cm}

\begin{lstlisting}[language=R, basicstyle=\tiny\ttfamily]
# Get weight matrix for multiple test samples
weights <- get_weight_matrix_distributed(model, X_test)

dim(weights$weight_matrix)  # m x n (test x train)
weights$sample_indices      # 1:n (all training indices)
weights$Y_values            # n x k (training Y values)
\end{lstlisting}

\vspace{0.3cm}

\textbf{Properties:}
\begin{itemize}
\item Each row sums to 1: $\sum_j W_{ij} = 1$
\item Sparse matrix (most weights are 0)
\item Prediction verification: $\hat{Y} = W \cdot Y_{\text{train}}$
\end{itemize}

\end{frame}

% ============================================================================
% SLIDE 11: Use Cases
% ============================================================================
\begin{frame}[fragile]{Practical Use Cases}

\begin{block}{Use Case 1: Find Influential Training Samples}
\end{block}

\begin{lstlisting}[language=R, basicstyle=\tiny\ttfamily]
weights <- get_weight_matrix_distributed(model, X_test)

# For first test sample, find top 5 influential training samples
top_5 <- order(weights$weight_matrix[1, ], decreasing = TRUE)[1:5]
cat("Most influential training samples:", top_5, "\n")
cat("Their weights:", weights$weight_matrix[1, top_5], "\n")
\end{lstlisting}

\vspace{0.3cm}

\begin{block}{Use Case 2: Compare Prediction Methods}
\end{block}

\begin{lstlisting}[language=R, basicstyle=\tiny\ttfamily]
# Mean-based RMSE
rmse_mean <- sqrt(mean((model$residuals$mean_based)^2))

# Parameter-based RMSE
rmse_param <- sqrt(mean((model$residuals$param_based)^2))

cat("RMSE (mean-based):", rmse_mean, "\n")
cat("RMSE (param-based):", rmse_param, "\n")
\end{lstlisting}

\end{frame}

% ============================================================================
% SLIDE 12: Performance Considerations
% ============================================================================
\begin{frame}{Performance & Scalability}

\begin{columns}[T]

\begin{column}{0.48\textwidth}
\textbf{Memory Usage:}

\begin{table}
\tiny
\begin{tabular}{lrr}
\toprule
\textbf{Mode} & \textbf{Memory} & \textbf{Factor} \\
\midrule
\texttt{store\_samples = FALSE} & 400 KB & 1× \\
\texttt{store\_samples = TRUE} & 20 MB & 50× \\
\bottomrule
\end{tabular}
\end{table}

\vspace{0.1cm}
\tiny{(1K samples, 100 trees, 5 components)}

\vspace{0.3cm}

\textbf{Prediction Speed:}

\begin{table}
\tiny
\begin{tabular}{lc}
\toprule
\textbf{Method} & \textbf{Speed} \\
\midrule
Pre-computed leaf & $\sim$1ms \\
Weight-based & $\sim$10ms \\
\bottomrule
\end{tabular}
\end{table}

\end{column}

\begin{column}{0.48\textwidth}
\textbf{Parallelization Speedup:}

\begin{itemize}
\item Sequential (1 core): baseline
\item 4 cores: $\sim$3.5× faster
\item 8 cores: $\sim$6-7× faster
\item Efficiency: 85-90\%
\end{itemize}

\vspace{0.3cm}

\textbf{Recommendations:}

\begin{itemize}
\item \textbf{Production:} \\
\texttt{store\_samples = FALSE}, many cores
\item \textbf{Research/Analysis:} \\
\texttt{store\_samples = TRUE}, weight analysis
\item \textbf{Large $n$ ($>$10K):} \\
Consider MoM over MLE
\item \textbf{Many trees ($>$200):} \\
Always use parallelization
\end{itemize}

\end{column}

\end{columns}

\end{frame}

% ============================================================================
% SLIDE 13: Complete Workflow
% ============================================================================
\begin{frame}[fragile]{Complete Workflow Example}

\begin{lstlisting}[language=R, basicstyle=\tiny\ttfamily]
library(DirichletForestParallel)
library(MCMCpack)

# 1. PREPARE DATA
n <- 500; p <- 4; k <- 3
X <- matrix(rnorm(n * p), n, p)
Y <- rdirichlet(n, alpha = c(2, 3, 4))
X_test <- matrix(rnorm(10 * p), 10, p)

# 2. BUILD MODEL (with interpretability enabled)
model <- DirichletForest_distributed(
  X, Y, 
  B = 100,                   # 100 trees
  d_max = 10,                # Max depth 10
  method = "mom",            # Method of moments
  store_samples = TRUE,      # Enable weight analysis
  n_cores = 4                # 4 parallel workers
)

# 3. EXAMINE FITTED VALUES
print(head(model$fitted$mean_based))      # In-sample predictions
print(head(model$residuals$mean_based))   # Residuals

# 4. MAKE PREDICTIONS (two modes)
pred_fast <- predict_distributed_forest(model, X_test, use_leaf_predictions = TRUE)
pred_weight <- predict_distributed_forest(model, X_test, use_leaf_predictions = FALSE)

# 5. ANALYZE SAMPLE WEIGHTS
weights <- get_weight_matrix_distributed(model, X_test)
top_influential <- order(weights$weight_matrix[1, ], decreasing = TRUE)[1:5]
print(top_influential)

# 6. CLEANUP (especially important on Windows)
cleanup_distributed_forest(model)
\end{lstlisting}

\end{frame}

% ============================================================================
% SLIDE 14: Summary
% ============================================================================
\begin{frame}{Summary: When to Use What}

\begin{table}
\small
\begin{tabular}{lp{5cm}p{4cm}}
\toprule
\textbf{Scenario} & \textbf{Recommended Settings} & \textbf{Why} \\
\midrule
Production deployment & 
\texttt{store\_samples = FALSE}, \newline \texttt{n\_cores = -1} & 
Fast, memory-efficient \\
\midrule
Research \& analysis & 
\texttt{store\_samples = TRUE}, \newline \texttt{use\_leaf\_predictions = FALSE} & 
Full interpretability \\
\midrule
Large dataset ($n > 10K$) & 
\texttt{method = "mom"}, \newline many cores & 
Computational efficiency \\
\midrule
Small dataset ($n < 500$) & 
\texttt{method = "mle"}, \newline \texttt{n\_cores = 1} & 
Better accuracy \\
\midrule
Uncertainty quantification & 
\texttt{store\_samples = TRUE}, \newline weight analysis & 
Access to distribution \\
\bottomrule
\end{tabular}
\end{table}

\vspace{0.3cm}

\begin{alertblock}{Key Takeaway}
\textbf{Trade-off:} Memory/Speed $\leftrightarrow$ Interpretability/Accuracy
\end{alertblock}

\end{frame}

% ============================================================================
% FINAL SLIDE
% ============================================================================
\begin{frame}{Thank You!}

\centering
\Large

\textbf{DirichletForestParallel}

\vspace{0.5cm}

\normalsize
High-performance compositional data regression with:
\begin{itemize}
\item Flexible storage modes
\item Parallel processing
\item Multiple prediction types
\item Full interpretability (optional)
\end{itemize}

\vspace{1cm}

Questions?

\end{frame}

\end{document}